% ============================================================
%  packages.tex — Configuración y estilo del CV
%  (No hace falta tocar este archivo para editar el contenido)
% ============================================================

% --- Página y tipografía --------------------------------------
\usepackage[a4paper, top=1.4cm, bottom=1.4cm, left=1.5cm, right=1.5cm]{geometry}
\usepackage[utf8]{inputenc}
\usepackage[T1]{fontenc}
\usepackage[spanish,es-noshorthands]{babel}
\usepackage[default]{lato}          % Sans-serif moderna y legible
\usepackage{parskip}                % Sin sangrías, espacio entre párrafos
\setlength{\parskip}{6pt}
\usepackage{microtype}

% --- Color, iconos y utilidades -------------------------------
\usepackage{xcolor}
\usepackage{graphicx}               % Foto
\usepackage{fontawesome5}
\usepackage{enumitem}
\usepackage{paracol}                % Layout a dos columnas
\usepackage{tikz}                   % Puntos de nivel de idiomas
\usepackage[colorlinks=true, urlcolor=accent, linkcolor=accent]{hyperref}

% --- Paleta ----------------------------------------------------
% Acento: azul petróleo — formal pero con carácter.
% Si quieres otro tono, cambia solo estas líneas.
\definecolor{accent}{HTML}{0E6BA8}      % Azul petróleo
\definecolor{accentdark}{HTML}{084B7A}  % Variante oscura (nombre, títulos)
\definecolor{accentlight}{HTML}{E3F0F8} % Fondo de las etiquetas de habilidades
\definecolor{body}{HTML}{3A3F45}        % Gris oscuro para el texto
\definecolor{soft}{HTML}{707880}        % Gris medio para fechas y detalles

\color{body}

% --- Cabecera ---------------------------------------------------
% \cvname{Nombre}  \cvrole{Titular profesional}
\newcommand{\cvname}[1]{{\fontsize{28}{32}\selectfont\bfseries\color{accentdark}#1\par}}
\newcommand{\cvrole}[1]{\vspace{2pt}{\large\color{accent}#1\par}}

% Foto circular con borde de acento: \cvphoto{ruta}
\newcommand{\cvphoto}[1]{%
  \begin{tikzpicture}
    \begin{scope}
      \clip (0,0) circle (1.6cm);
      \node at (0,-0.12) {\includegraphics[width=3.4cm]{#1}};
    \end{scope}
    \draw[accent, line width=1.6pt] (0,0) circle (1.6cm);
  \end{tikzpicture}%
}

% Línea de contacto: \cvcontact{icono}{texto}
\newcommand{\cvcontact}[2]{{\color{accent}#1}\hspace{5pt}{\small #2}}
\newcommand{\cvsep}{\hspace{9pt}{\color{soft}\textbar}\hspace{9pt}}

% --- Títulos de sección -----------------------------------------
\newcommand{\cvsection}[1]{%
  \vspace{8pt}%
  {\large\bfseries\color{accentdark}\MakeUppercase{#1}}\par\nopagebreak
  \vspace{\dimexpr 1pt-\parskip\relax}%
  {\color{accent}\rule{28pt}{2.4pt}}{\color{accentlight}\rule{\dimexpr\linewidth-28pt\relax}{2.4pt}}\par\nopagebreak
  \vspace{\dimexpr 2pt-\parskip\relax}%
}

% --- Entradas (formación / experiencia) --------------------------
% \cventry{Fechas}{Título}{Lugar / institución}{Descripción (opcional)}
\newcommand{\cventry}[4]{%
  {\bfseries #2}\hfill{\small\color{soft}#1}\par\nopagebreak
  \vspace{\dimexpr 1pt-\parskip\relax}%
  {\color{accent}\small #3}\par
  \ifx\relax#4\relax\else
    \vspace{\dimexpr 1pt-\parskip\relax}%
    {\small #4\par}%
  \fi
  \vspace{3pt}%
}

% --- Etiquetas de habilidades ("pastillas") ----------------------
% Envuelve los grupos de \cvtag en el entorno cvtags para que las
% etiquetas salten de línea sin salirse del margen.
\newenvironment{cvtags}{\raggedright}{\par}
\newcommand{\cvtag}[1]{%
  \tikz[baseline=(txt.base)]{
    \node[fill=accentlight, text=accentdark, rounded corners=3pt,
          inner xsep=7pt, inner ysep=4pt, font=\small] (txt) {#1};
  }\hspace{1pt}\vspace{3pt}
}

% --- Nivel de idioma con puntos ----------------------------------
% \cvlang{Idioma}{nivel de 0 a 5}
\newcommand{\cvlang}[2]{%
  #1\hfill
  \tikz[baseline=-2.5pt]{
    \foreach \i in {1,...,5}{
      \ifnum\i>#2
        \fill[accentlight] (\i*0.3, 0) circle (0.085);
      \else
        \fill[accent] (\i*0.3, 0) circle (0.085);
      \fi
    }
  }\par\vspace{4pt}%
}

% --- Listas compactas --------------------------------------------
\setlist[itemize]{leftmargin=12pt, itemsep=2pt, topsep=2pt, label={\color{accent}\small\textbullet}}

\pagestyle{empty}
